\section{超图关联及完整性检查表}
\noindent 对实际有限转录，以每个标量决策为节点，每条约束的真实非零依赖集为超边；分别保存变量块表达，不能混用两种粒度。[R6]
\begin{align}
 n&=|\mathcal X|,\quad m=|\mathcal C|,\quad
 N_I=\sum_{c\in\mathcal C}|\operatorname{supp}F_c|,\\
 r&=\max_c|\operatorname{supp}F_c|,\quad
 d_{max}=\max_x|\{c:x\in\operatorname{supp}F_c\}|.
 \label{eng:eq:hypergraph}
\end{align}
\noindent 目标函数的依赖另外登记；辅助变量和消元可改变上述计数，必须固定表示再比较。天气是外生输入，多个航次读取同一风值并不自动构成资源超边。半径和度数也不直接证明超图算法更优。
\begin{center}\small
\begin{tabular}{p{.24\linewidth}p{.42\linewidth}p{.24\linewidth}}\toprule
约束族 & 真实耦合对象 & 排除的问题\\\midrule
整塔与任务流 & 同塔9项、航次、顺序 & 部分照片冒充整塔；孤立环\\
连续飞行与相机 & 航迹、时间、姿态、云台、焦距 & 瞬移、模糊、视野越界\\
道路与会合 & 路段、连续车体、停靠、飞行端点 & 无路可达却起降；车辆瞬移\\
时空冲突 & 同时刻车辆/飞机/操作区 & 只检查空间或只检查排程\\
电池与充电 & 航次组成、物理身份、电量和通道 & 重复使用、空中充电、免费补能\\
天气与能耗 & 时空轨迹、空速、功率、窗口 & 全日套首小时；使用未来真值\\
重规划与终态 & 已验收、承诺、当前电量、返库 & 重置状态、遗失空中飞机\\\bottomrule
\end{tabular}\end{center}
\noindent 本模型的$n,m,N_I,r,d_{max}$要在237塔完整转录后实测，当前均为未知；不能用2133任务数、19类旧变量块或64维网络隐藏向量替代。当前原始源数据、修正层、变量域、约束索引及缺失输入在JSON中绑定。

\noindent 对比超图方法与一般方法时，二者必须使用相同参考几何修正、完整任务/道路域、天气信息截止、物理约束、初末状态和目标。已有9任务与其他公开案例不是本237塔联合模型的求解结果。
\clearpage

